% THIS IS A LATEX TEMPLATE FILE FOR PAPERS INCLUDED IN THE
% *Anthology of Computers and the Humanities*. ADD THE OPTION
% 'final' WHEN CREATING THE FINAL VERSION OF THE PAPER. 
% DO NOT change the documentclass
%\documentclass[final]{anthology-ch} % for the final version
\documentclass{anthology-ch}         % for the submission

% LOAD LaTeX PACKAGES
\usepackage{booktabs}
\usepackage{graphicx}
% ADD your own packages using \usepackage{}
\usepackage{minted}
\usemintedstyle{manni}
\usepackage{xcolor}
\definecolor{LightGray}{gray}{0.95}
\setminted[python]{bgcolor=LightGray}
\usepackage{relsize}

% TITLE OF THE SUBMISSION
% Change this to the name of your submission
\title{Undate in Action}

% AUTHOR AND AFFILIATION INFORMATION
% For each author, include a new call to the \author command, with
% the numbers in brackets indicating the associated affiliations 
% (next section) and ORCID-ID for each author.  
\author[1]{Rebecca Sutton Koeser}[
  orcid=0000-0002-8762-8057
]
% There should be one call to \affiliation for each affiliation of
% the authors. Multiple affiliations can be given to each author
% and an affiliation can be given to multiple authors. 
\affiliation{1}{Center for Digital Humanities, Princeton University, Princeton, New Jersey, USA}

% KEYWORDS
% Provide one or more keywords or key phrases seperated by commas
% using the following command
\keywords{computers, humanities proceedings}

% METADATA FOR THE PUBLICATION
% This will be filled in when the document is published; the values can
% be kept as their defaults when the file is submitted
\pubyear{2025}
\pubvolume{1}
\pagestart{1}
\pageend{1}
\conferencename{Proceedings of Conference XXX}
\conferenceeditors{Editor1 Editor2}
\doi{00000/00000}  

\addbibresource{bibliography.bib}

%%%%%%%%%%%%%%%%%%%%%%%%%%%%%%%%%%%%%%%%%%%%%%%%%%%%%%%%%%%%%%%%%%%%%%%%%%%
% HERE IS THE START OF THE TEXT
\begin{document}

\maketitle

\begin{abstract}
This LaTeX template helps you typeset and format a paper for the Computational Humanities Research conference in the ACH Anthology. This template helps you adhere to the the required specifications and provides an example of how your paper should look. In practice, the abstract of the paper here should be a one-paragraph summary of the outline and main contributions of the paper. 
\end{abstract}

\section{Introduction} 


\texttt{undate} is a beta Python library for working with uncertain or partially known dates, with support for multiple calendars.

https://github.com/dh-tech/undate-python/

Originally started at a DHtech hackathon! Builds on experiences from
other digital humanities projects.

\subsection{How is \texttt{Undate} like \texttt{datetime.date} ?} \label{sec:intro_details}


Both can be initialized by specifying numeric values for year, month, and day. Both can be printed using the default serialization (ISO8601, or YYYY-MM-DD), and we can compare them.

\begin{minted}{python}
import datetime

from undate import Undate

year, month, day = 2000, 11, 7

# these are equivalent and we initialize them the same way
dt_november7 = datetime.date(year, month, day)
november7 = Undate(year, month, day)

assert november7 == dt_november
\end{minted}


\subsection{How is \texttt{Undate} \textit{not} like \texttt{datetime.date}?}

An \texttt{Undate} can be initialized with partial information.  \texttt{Undate} can also take an optional label, since sometimes it's useful to name a date.  It can be initialized with year and month, just year, or even month and day with no year.

\begin{minted}{python}
# November 2000
november = Undate(year, month, label="November 2000")
# Year 2000
year2k = Undate(year, label="Y2K")
# November 7 in an unknown year
november7_some_year = Undate(month=month, day=day, label="Some November 7")
november7.label = "November 7, 2000"

# sometimes names are important
easter1916 = Undate(1916, 4, 23, label="Easter 1916")
\end{minted}


If you try to do that with \texttt{datetime.date}, you get a
\texttt{TypeError} because all three fields are required.

\begin{minted}{python}
try:
    datetime.date(year, month)
except TypeError as err:
    print(err)

"function missing required argument 'day' (pos 3)"
\end{minted}


Each of these \texttt{Undate} objects can be printed out in a standard
format; the \texttt{Undate} class also tracks how precisely a date is
specified, and can calculate the duration.  If we iterate over the
example partially-known dates initialized in the code block above and print label, formatted date, precision, and duration in days, we get output like this (here and throughout, formatting added for readability):

\begin{minted}{python}
for example_date in [
    november,
    year2k,
    november7_some_year,
    november7,
    easter1916,
]:
print(f"""{example_date.label} {example_date}
- Date precision: {example_date.precision}
- Duration in days: {example_date.duration().days}""")
\end{minted}

\paragraph{November 2000} \texttt{2000-11}
\begin{itemize}
    \item Date precision: MONTH
    \item Duration in days: 30
\end{itemize}

\paragraph{Y2K} \texttt{2000}
\begin{itemize}
    \item Date precision: YEAR
    \item Duration in days: 366
\end{itemize}

\paragraph{Some November 7} \texttt{-}\texttt{-11-07}
\begin{itemize}
   \item  Date precision: DAY
   \item Duration in days: 1
\end{itemize}

\paragraph{November 7, 2000} \texttt{2000-11-07}
\begin{itemize}
    \item Date precision: DAY
    \item Duration in days: 1
\end{itemize}

\paragraph{Easter 1916} \texttt{1916-04-23}
\begin{itemize}
    \item Date precision: DAY
    \item Duration in days: 1
\end{itemize}

We can also do some simple calculations, like checking whether one date falls within another date.
\begin{minted}{python}
assert november in year2k   # 2000-11 in 2000
assert november7 in november  # 2000-11-07 in 2000-11
assert easter1916 not in year2k  # 1916-02-23 not in 2000
assert november7_some_year not in year2k  # --11-07 not in 2000
\end{minted}

\section{What else can \texttt{Undate} do?}

\subsection{Partially known values}

In the last set of examples, the values used to initialize \texttt{Undate} instances were all integers and parts of the date were either known or unknown. But what if you know part of a
date?

You can initialize \texttt{Undate} with strings and use \texttt{X} to
indicate unknown values.  (We chose \texttt{X} based on the \href{https://www.loc.gov/standards/datetime/}{Extended Date Time Format (EDTF)} specification).

Here we initialize two dates with partially known information, where a year or month cannot be precisely pinned down but some information is known. As before, we can iterate over these dates
and output label, formatted date, precision, and duration:

\begin{minted}{python}
someyear_1900s = Undate("19XX", label="1900s")
late2022 = Undate(2022, "1X", label="late 2022")

for example_date in [someyear_1900s, late2022]:
    print(f"""{example_date.label} {example_date}
- Date precision: {example_date.precision}
- Duration in days: {max(example_date.duration().days)}
)
\end{minted}

\paragraph{1900s} \texttt{19XX}
\begin{itemize}
    \item Date precision: YEAR
    \item Duration in days: 366
\end{itemize}
    
\paragraph{late 2022} \texttt{2022-1X}
\begin{itemize}
    \item Date precision: MONTH
    \item Duration in days: 31
\end{itemize}    

    
If you try that with \texttt{datetime.date} you get another \texttt{TypeError}.    

\begin{minted}{python}
try:
    datetime.date("19XX", 1, 1)
except TypeError as err:
    print(err)

"'str' object cannot be interpreted as an integer"
\end{minted}

\subsection{Uncertain durations}

If you were paying close attention, you may have noticed I used the
\texttt{max()} function in the code block above where I output the duration for the partially known dates.

The most recent version of \texttt{undate} (version 0.5) includes experimental support for 
uncertain time deltas.

\begin{minted}{python}
print(f"""{someyear_1900s.label} {someyear_1900s}
- duration: {someyear_1900s.duration()}
- duration in days: {someyear_1900s.duration().days}

{late2022.label} {late2022}
- duration: {late2022.duration()}
- duration in days: {late2022.duration().days}
"""
\end{minted}

\paragraph{1900s} \texttt{19XX}
\begin{itemize}
    \item duration: \texttt{UnDelta(days=[365,366])}
    \item duration in days: \texttt{UnInt(lower=365, upper=366)}
\end{itemize}

\paragraph{late 2022} \texttt{2022-1X}
\begin{itemize}
    \item duration: \texttt{UnDelta(days=[30,31])}
    \item duration in days: \texttt{UnInt(lower=30, upper=31)}
\end{itemize}
    
    
Even without precise information, we can still do some useful comparisons. February of an unknown year is still shorter than October, November, or December.

\begin{minted}{python}
some_february = Undate(month=2, label="February of some year")
assert some_february.duration() < late2022.duration()

# February of an unknown year is shorter than 
# an uncertain month 1X (October, November, or December) 
some_february.duration() < late2022.duration() = True

assert some_february.duration() < late2022.duration()

print(f"""### {some_february.label} {some_february}
- duration: {some_february.duration()}
- duration in days: {some_february.duration().days}
""")

\end{minted}

\paragraph{February of some year} \texttt{-}{-02}
\begin{itemize}
    \item duration: \texttt{UnDelta(days=[28,29])}
    \item duration in days: \texttt{UnInt(lower=28, upper=29)}
\end{itemize}
    
\section{Example use cases from specific projects}\label{example-use-cases-from-specific-projects}


\textbf{TODO}: briefly describe here how undate generalizes work from projects and is informed by project-specific context, and we hope to keep that way in future (quote sw article?)

\subsection{Princeton Geniza Project}\label{princeton-geniza-project}

The \href{https://geniza.princeton.edu/}{\textit{Princeton Geniza Project}} (PGP) is
a long-running project focused on materials from a synagogue in Cairo; these contents are
primarily medieval, and largely written in Hebrew script. Because they are older and fragmentary, many of them cannot be dated;  those documents that do have dates use a variety of calendars — the Hebrew \textit{Anno Mundi} and Seleucid calendars, as well as the
Islamic \textit{Hijri} calendar. 

An \texttt{Undate} instance preserves any initial values for year, month, and day in the original calendar. As you see in Table \ref{tab:pgp-parsed-dates}, the ISO format \texttt{Undate} corresponds to the original date rather than the standardized date.

\begin{table}[h]
  \centering 
  \begin{tabular}{rlllllll}
    \toprule
Original Date & Calendar & Standard date & \texttt{Undate} & Date precision & Weekday \\ 
\midrule\addlinespace[2.5pt]
1570 & Seleucid & 1259 & 1570 & year &  \\
Tammuz 1288 & Seleucid & 0977-06-21/0977-07-19 & 1288-04 & month &  \\
19 Adar 1427 & Seleucid & 1116-03-05 & 1427-12-19 & day & Sunday \\
\midrule
4890 & Anno Mundi & 1129-09-16/1130-09-05 & 4890 & year &  \\
Av 5564 & Anno Mundi & 1804-07-09/1804-08-07 & 5564-05 & month &  \\
10 Nisan 4716 & Anno Mundi & 0956-03-24 & 4716-01-10 & day & Monday \\
\midrule
537 & Hijrī & 1142-07-27/1143-07-15 & 0537 & year &  \\
Shawwāl 425 & Hijrī & 1034-08-29/1034-09-07 & 0425-10 & month &  \\
5 Safar 584 & Hijrī & 1188-04-05 & 0584-02-05 & day & Tuesday \\
\bottomrule
    \bottomrule
  \end{tabular}
  \caption{A sampling of dates from PGP documents in different calendars; dates have been parsed by \texttt{Undate}, which was then used to report precision and weekday.}
  \label{tab:pgp-parsed-dates}
\end{table}

Under the hood, \texttt{undate} calculates the earliest and latest possible dates in range for comparison and sorting. When an \texttt{Undate} object is initialized with a different calendar, the earliest and latest dates are based on conversion to the Gregorian calendar so that dates can be compared and used together across calendars. (For the code used to parse PGP document dates displayed in this example, refer to Appendix \ref{appdx:parse-pgp-dates}). When working across calendars, dates with the same precision (year, month) often have different durations in days.

\begin{table}[h]
  \centering 
  \begin{tabular}{rlllllll}
    \toprule
Original Date & Calendar & \texttt{Undate} & Earliest & Latest & Date Precision & Duration \\ 
\midrule
1570 & Seleucid & 1570 & 1258-09-07 & 1259-09-26 & year & 383 \\
1474 & Seleucid & 1474 & 1162-09-18 & 1163-09-06 & year & 354 \\
Tevet 1363 & Seleucid & 1363-10 & 1051-12-14 & 1052-01-11 & month & 29 \\
{[12]91} & Seleucid & 1291 & 0979-09-30 & 0980-09-17 & year & 353 \\
Heshvan 1453 & Seleucid & 1453-08 & 1141-10-11 & 1141-11-08 & month & UnInt(29,30) \\
1347 & Seleucid & 1347 & 1035-09-12 & 1036-09-28 & year &  383 \\
\midrule
Shevat (4)791 & Anno Mundi & 4791-11 & 1031-01-03 & 1031-02-01 & month & UnInt(29,30) \\
4812 & Anno Mundi & 4812 & 1051-09-15 & 1052-09-03 & year & 385 \\
\midrule
502 & Hijrī & 0502 & 1108-08-18 & 1109-08-06 & year & 354 \\
Muḥarram 557 & Hijrī & 0557-01 & 1161-12-28 & 1162-01-26 & month & 30 \\
Dhū l-Qa'da 525 & Hijrī & 0525-11 & 1131-10-02 & 1131-10-31 & month & 29 \\
\bottomrule
    \bottomrule
  \end{tabular}
  \caption{Another sampling of dates from PGP documents in different calendars parsed by \texttt{Undate}. Earliest and latest dates are calculated in Gregorian / proleptic Gregorian calendar. Semantic durations like year and month result in different durations in days when working with multiple calendars.}
  \label{tab:pgp-date-durations}
\end{table}

%% should any of these be UnInt? might be a bug in undate...

An \texttt{Undate} object is aware of date precision, which means that in a dataset with mixed precision dates like PGP, once the dates are parsed we can easily see the variation (Table \ref{tab:pgp-parsed-dates}).\footnote{With the caveat that in this case the results are somewhat skewed by the dates we're able to parse; PGP data includes the modifiers not yet supported by \texttt{undate}, which were ignored during parsing}.

\begin{table}[h]
  \centering 
  \begin{tabular}{cc}
    \toprule
    Original Date precision & Documents \\ 
    \midrule
    year & 831 \\
    month & 1,021 \\
    day & 1,558 \\
    \bottomrule
  \end{tabular}
  \caption{Totals for PGP documents and date precision}
    \medskip
    \small
    Only includes documents with standard dates that can be parsed by \texttt{undate}.
  \label{tab:example}
\end{table}


todo: add doc frequency by weekday from undate sw paper


\subsection{Shakespeare and Company
Project}\label{shakespeare-and-company-project}

\href{https://shakespeareandco.princeton.edu/}{\textit{Shakespeare and Company
Project}} is based on the materials from Sylvia Beach's famous
English-language lending library which operated in Paris in the 1920s
and 1930s.

This project includes borrowing events with unknown years; with
\texttt{undate} we can calculate how long a book was borrowed even when
we don't know the exact year.

First we define a couple of methods that we'll use with the data; one to
parse the date and calculate duration in days, and another for an
indicator whether the year was known or not.


\begin{table}[h]
  \centering 
  \begin{tabular}{lllrllrrl}
    \toprule
    Start Date & End Date & Duration & Title & Author \\ 
    \midrule
--01-19 & --01-24 & 5 & \textit{A Hazard of New Fortunes} & Howells, William Dean \\
--01-24 & --05-30 & 126 & \textit{Boston Cooking-School Cook Book} & Farmer, Fannie Merritt \\
--01-24 & --03-20 & 55 & \textit{The Landlord at Lion's Head} & Howells, William Dean \\
--01-24 & --05-30 & 126 & \textit{Mrs. Rorer's New Cook Book} & Rorer, Sarah Tyson \\
--01-24 & --03-20 & 55 & \textit{Indian Summer} & Howells, William Dean \\
--01-24 & --03-20 & 55 & \textit{A Foregone Conclusion} & Howells, William Dean \\
--01-24 & --03-20 & 55 & \textit{Dr. Breen's Practice} & Howells, William Dean \\
1920-04-29 & 1920-06-03 & 35 & \textit{The Letters of George Meredith} & Meredith, George \\
    \bottomrule
  \end{tabular}
  \caption{Example table and table caption.}
  \label{tab:example}
\end{table}




-------


\paragraph{Another tip.} In some cases, it may be helpful to use \texttt{paragraph} to title individual paragraphs. For example, if a section describes features for a classifier, you can optionally title each paragraph with the name of each feature. 

\section{Elements}

\subsection{Citing elements}

Here are some examples of how to construct and reference common elements in LaTeX. References to elements such as tables, figures, equations and sections make use of \texttt{label} names that you set. References to citations should use the labels you indicate in \texttt{bibliography.bib}. Change all of these examples and values with your own data. 

We can cite Table~\ref{tab:example} as well as Figure~\ref{fig:example}, and we also cite an example paper \cite{tettoni2024discoverability}.
We can also include mathematical notations, such as:
\begin{align}
f(y) &= x^2. \label{fig:squared}
\end{align}
The line number of the equation can be cited as
Equation~\ref{fig:squared}. You can also cite multiple papers together \cite{barré2024latent, levenson2024textual, bambaci2024steps}, and reference figures or tables indirectly in parentheses (Figure~\ref{fig:example_bigger}). You can also cite other sections or subsections of your paper, such as \S\ref{sec:intro_details}. 


\begin{table}[h]
  \centering 
  \begin{tabular}{cc}
    \toprule
    Column Name 1 & Column Name 2\\
    \midrule
    d1 & d2 \\
    d1 & d2 \\
    d1 & d2 \\
    \bottomrule
  \end{tabular}
  \caption{Example table and table caption.}
  \label{tab:example}
\end{table}


\subsection{Required specifications}

Tables and figures should \textit{not} appear at the top of the first page above the paper title and abstract, but can be placed within the main text, as exemplified by Table~\ref{tab:example}. They may also be placed at the top of non-first pages, as exemplified by Figures~\ref{fig:example} and \ref{fig:example_bigger}. Figures and tables discussed in the main text should appear \textit{before} the References section. Supplementary materials should be referenced by their relevant Appendix section, such as Appendix~\ref{appdx:first}. 

Do \textit{not} change the font size of table and figure captions, or the spacing between text lines, section/subsection titles, tables, figures, and captions. You should size your figures and tables so that they stay within the \texttt{linewidth} of the paper. 

\begin{figure}[t!]
  \centering
  \includegraphics[width=0.4\linewidth]{640x480.png}
  \caption{Example figure and figure caption.}
  \label{fig:example}
\end{figure}

\begin{figure}[t!]
  \centering
  \includegraphics[width=0.4\linewidth]{640x480.png}
  \includegraphics[width=0.4\linewidth]{640x480.png}
  \caption{Example figure, where two \texttt{.png} are placed side by side.}
  \label{fig:example_bigger}
\end{figure}

%\section*{Acknowledgements}
%
%This unnumbered section should be blank when submitting your paper. After review, you may include lists of people and organizations who supported the work.

% Print the biblography at the end. Keep this line after the main text of your paper, and before an appendix. 
\printbibliography

% You can include an appendix using the following command
\appendix

\section{Parsing PGP document dates in multiple calendars} \label{appdx:parse-pgp-dates}

Princeton Geniza Project data includes documents with original dates in multiple calendars.  \texttt{undate} does not yet provide an omnibus parser, so dates must be parsed according to their a known, supported calendar.  The dates in this dataset are complicated and messy, and include qualifiers and uncertainty not yet supported by \texttt{undat}e parsing. For demonstration purposes, this example code removes and ignores those complexities.

Adapted from an \href{https://github.com/dh-tech/undate-python/blob/main/examples/pgp_dates.ipynb}{example notebook} included in the \texttt{undate} software repository.

\begin{minted}{python}
import re

from lark.exceptions import UnexpectedEOF
from lark.visitors import VisitError
import polars as pl


# set this to True to see details about parsing
VERBOSE_PARSE_OUTPUT = False


def remove_ordinals(val):
    # utility method to turn ordinals into numerals
    # i.e., 1st, 2nd, 3rd -> 1, 2, 3
    return re.sub(r"(\d+)(st|nd|rd|th)", "\\1", val)


def parse_original_date(doc_date_original, doc_date_calendar):
    # Map PGP original calendar to undate calendar converter
    undate_calendar = None
    if doc_date_calendar == "Anno Mundi":
        undate_calendar = "Hebrew"
    elif doc_date_calendar == "Hijrī":
        undate_calendar = "Islamic"
    elif doc_date_calendar == "Seleucid":
        undate_calendar = "Seleucid"

    if undate_calendar:
        value = doc_date_original

        # some dates have unknown digits, e.g. 1[.] Kislev 48[..] or 152[.]
        # Replace with zeros for now, since the calendar parsers
        # doesn't yet support this syntax
        # (even though Undate supports unknown digits)
        if "[." in value:
            if VERBOSE_PARSE_OUTPUT:
                print(f"ignoring missing digits for now {value}")
            value = (
                value.replace("[.]", "0")
                .replace("[..]", "00")
                .replace("[...]", "000")
            )

        # some dates have inferred numbers, 
        # e.g. Friday, [25] Nisan [4810] or 8 Elul (4)811'
        # for now, we strip out brackets before parsing;
        # in future, we could indicate uncertainty based on these
        value = (
            value.replace("[", "")
            .replace("]", "")
            .replace("(", "")
            .replace(")", "")
        )

        # for now, remove modifiers not supported by the undate parser:
        #   Late Tevet 4903, Last decade of Kislev 5004, first third of ...
        modifiers = [
            "Late ",
            "(first|middle|last)( third|half|decade|tenth)? (of )?",
            "(Beginning|end) of ",
            "last day",
            "First 10 days",
            " of",
            "spring",
            "decade ",
            "night, ",
        ]
        for mod in modifiers:
            value = re.sub(mod, "", value, flags=re.I)

        # about 62 have ordinals; strip them out
        value = remove_ordinals(value)

        # parse the simplified, cleaned up date string 
        # with the specified calendar
        try:
            return Undate.parse(value, undate_calendar)
        except (VisitError, ValueError, UnexpectedEOF) as err:
            if VERBOSE_PARSE_OUTPUT:
                print(f"Parse error : {value} ({undate_calendar}): {err}")
            return None

# Use Polars to load data file into a DataFrame
# Filter to documents with standard dates
pgp_documents_df = pl.read_csv("pgp_documents.csv") \     
    .filter(
        pl.col("doc_date_standard").is_not_null()
    )
    
# Create a struct of original date and calendar,
# which we can pass to our parse method.
# If parsing succeeds, returns and Undate or UndateInterval
pgp_documents_df = pgp_documents_df.with_columns(
    undate_orig=pl.struct(
        "doc_date_original", "doc_date_calendar"
    ).map_elements(
        lambda row: parse_original_date(
            row["doc_date_original"], row["doc_date_calendar"]
        ),
        return_dtype=pl.datatypes.Object,
    )
)
days = [
    "Monday",
    "Tuesday",
    "Wednesday",
    "Thursday",
    "Friday",
    "Saturday",
    "Sunday",
]

# filter to records that were successfully parsed; determine weekday for use with day-level precision dates

pgp_documents_df = (
    pgp_documents_df.filter(pl.col("undate_orig").is_not_null())
    .with_columns(
        pl.col("type").fill_null("Unknown"),  # set null type to unknown
        orig_date_precision=pl.col("undate_orig").map_elements(
            lambda x: str(x.precision).lower(),
            return_dtype=pl.datatypes.String,
        ),
        undate_weekday=pl.col("undate_orig").map_elements(
            lambda x: x.earliest.weekday, return_dtype=pl.datatypes.Int32
        ),
    )
    .with_columns(
        undate_weekday_name=pl.col("undate_weekday").map_elements(
            lambda x: days[x], return_dtype=pl.datatypes.String
        )
    )
)

# filter to relevant fields and display the first ten rows
pgp_documents_df.select(
    "pgpid",
    "type",
    "doc_date_original",
    "doc_date_calendar",
    "doc_date_standard",
    "undate_orig",
    "orig_date_precision",
    "undate_weekday_name",
).head(10)    
\end{minted}


\end{document}

\section{First Appendix Section} \label{appdx:first}

Appendix sections should be ordered by letters rather than numbers, and their contents do not count towards the paper's length limit. Appendix sections may also contain additional tables and figures.  

\end{document}
